\section{Choosing the right description}

This chapter introduced many descriptions of a line. At first this may seem like too much notation. But the purpose is not to collect formulas. The purpose is to learn how to choose a description that fits the problem.

The same line can be written in several ways. The best form depends on what is known and what is being asked.

\subsection*{When a point and direction are given}

If a problem gives a point and a direction vector, the parametric form is usually the most natural:
\[
P=A+t v.
\]
It keeps the given information visible. It also makes it easy to generate points and to understand movement along the line.

This form is also the safest first choice in space.

\subsection*{When two points are given}

If a problem gives two distinct points \(A\) and \(B\), first form the direction vector
\[
B-A.
\]
Then write
\[
P=A+t(B-A).
\]
From there, if needed, convert to another form.

The important point is that the two given points are not just coordinates to insert into a memorized formula. They define a displacement, and that displacement defines a direction.

\subsection*{When a normal vector is visible}

If a problem gives an equation
\[
Ax+By+C=0,
\]
then the normal vector
\[
n=(A,B)
\]
is immediately visible. This form is useful for distance questions, side-of-line questions, and comparisons using normal vectors.

If a point \(P_0\) and a normal vector \(n\) are given, use
\[
n\cdot(P-P_0)=0.
\]
This is often more meaningful than jumping directly to expanded coordinates.

\subsection*{When slope is useful}

If the line is non-vertical and slope is central to the question, then slope forms may be convenient:
\[
y-y_0=m(x-x_0)
\]
or
\[
y=mx+b.
\]
But these forms should be used with awareness. Vertical lines require a different description, such as
\[
x=c,
\]
or a parametric or general equation.

\subsection*{Which form helps which question?}

The following guide is often useful.

\begin{center}
\begin{tabular}{ll}
\toprule
Question & Useful description \\
\midrule
Generate points on the line & parametric form \\
Line through two points & two-point form, then parametric form \\
Check whether a point lies on the line & implicit or parametric form \\
Find an intersection in the plane & general form as a system \\
Check parallelism & direction vectors or normal vectors \\
Check perpendicularity & dot product of directions or normals \\
Compute distance from a point & general form \(Ax+By+C=0\) \\
Work in space & parametric form \\
\bottomrule
\end{tabular}
\end{center}

This table should not be used mechanically. It is a guide for choosing what to read from the line.

\subsection*{A complete reading example}

Consider the line
\[
\ell: 2x-3y+6=0.
\]
From this form we immediately read a normal vector:
\[
n=(2,-3).
\]
A direction vector can be chosen perpendicular to \(n\):
\[
v=(3,2),
\]
because
\[
(2,-3)\cdot(3,2)=6-6=0.
\]
To find a point on the line, set \(x=0\). Then
\[
-3y+6=0,
\]
so
\[
y=2.
\]
Thus \((0,2)\) lies on the line. A parametric form is
\[
P=(0,2)+t(3,2).
\]
The same line can therefore be read in two ways:
\[
2x-3y+6=0
\]
shows the normal vector and is good for distance or side tests, while
\[
P=(0,2)+t(3,2)
\]
shows a point and direction and is good for motion along the line.

\begin{remarkbox}
A good solution often changes the description of a line. This is not extra work. It is how we choose the language in which the problem becomes clear.
\end{remarkbox}

\subsection*{What to remember}

A line is a set of points. Equations and parametrizations are descriptions of that set. The main descriptions are not competitors; they are tools.

\begin{summarybox}
When working with lines, ask:
\begin{itemize}
\item Is the line in the plane or in space?
\item Am I given a point, two points, a direction vector, or a normal vector?
\item Which description makes the given information visible?
\item Which description makes the required property easiest to study?
\item Am I studying membership, intersection, parallelism, perpendicularity, or distance?
\item After computing, can I explain what the result means geometrically?
\end{itemize}
\end{summarybox}

The most important lesson is not that lines have many formulas. The most important lesson is that each formula reveals a different aspect of the same geometric object.